\chapter{Conclusions}

Effective Field Theories (EFTs) provide the most general and systematic way to describe a physical
system at energies below the characteristic scale of a more fundamental theory, encoding the
indirect effects of the heavy degrees of freedom in the coefficients of higher-dimensional
operators. Among the many EFTs developed in particle physics, those describing
beyond-the-Standard-Model effects are of particular interest. The Standard Model Effective Field
Theory (SMEFT) plays a key role, being the most general local, Lorentz-invariant effective theory
built from the Standard Model field content and invariant under its gauge group, with the Higgs in a
linearly realized $SU(2)_L$ doublet. Another widely studied example is the EFT of axion-like
particles (ALPs), which introduces a light pseudoscalar field whose interactions are suppressed by a
high-energy scale. Such fields naturally arise in theories with spontaneously broken approximate
global symmetries, where the ALP plays the role of a pseudo-Nambu--Goldstone boson, and they provide
a compelling framework to address several open questions in particle physics, such as the strong CP
problem, the flavor puzzle, the nature of dark matter, and the electroweak-scale naturalness
problem.

Within the EFT framework, the Wilson coefficients of higher-dimensional operators are
scale-dependent couplings that encode heavy-physics effects, and fundamental principles impose
non-trivial constraints on their allowed values. Therefore, to connect experimental data with the
underlying high-energy physics, it is essential to determine:
\begin{itemize}
    \item How Wilson coefficients evolve with energy, which is governed by the anomalous dimensions
    of the corresponding operators;
    \item Which regions of parameter space are consistent with fundamental principles, such as
    unitarity, causality, and locality of quantum field theory.
\end{itemize}

However, EFT computations are often technically demanding, since increasing operator dimension,
multiple operator insertions, and the presence of particles with spin quickly make Feynman-diagram
calculations intricate, while the number of diagrams grows factorially with the number of loops and
external legs. Moreover, Lagrangians carry gauge and off-shell redundancies, so individual diagrams
depend on conventions that cancel only in their sum.

To address the points above and overcome these challenges, in this thesis we employed on-shell
amplitude methods, a modern reformulation of perturbation theory that replaces the traditional
Feynman-diagram expansion with physical, gauge-invariant building blocks. In this framework,
scattering amplitudes are reconstructed by exploiting unitarity, locality, and factorization, and
are expressed in spinor-helicity variables, which reorganize the external kinematic data so that
gauge invariance and Lorentz covariance are manifest, yielding compact and transparent expressions.
This framework both streamlines calculations and reveals deep analytic and symmetry structures
underlying EFTs.

\subsubsection*{Anomalous Dimensions}

Predictions for physical processes are obtained by evaluating matrix elements of the EFT Lagrangian
at the energies probed by experiments. Therefore, the Lagrangian defined at the scale $\Lambda$ has
to be evolved down to the experimental scale $E \ll \Lambda$, which is governed by the anomalous
dimensions of the higher-dimensional operators. These control both the multiplicative
renormalization of the operators and their mixing, which determines how an experimental bound on one
operator affects the Wilson coefficients of the others. Moreover, anomalous-dimension matrices are
much sparser than naive power counting would suggest. In the diagrammatic approach, these zeros
emerge only after all diagrams are summed and the off-shell Green's basis is reduced to the physical
one, where they appear as accidental cancellations. Instead, in the on-shell framework they follow
from selection rules based on helicity, operator length, and angular momentum. In this context,
the contributions of this thesis are the following.
\begin{itemize}
    \item We generalized the on-shell extraction of anomalous dimensions from phase-space integrals
    of lower-loop amplitudes to the most general mixing among operators of different dimensions,
    including multiple operator insertions, up to two loops. Moreover, we included leading mass
    effects through the Higgs low-energy theorem, while keeping the formalism massless, which
    preserves both the simplicity of the massless spinor-helicity variables and the selection rules.
    These two aspects are essential
    for many EFTs and associated phenomenological studies. Indeed, in the presence of relevant and marginal
    couplings, operators of different dimensions mix under renormalization. Furthermore,
    chirality-flipping observables, such as electric and magnetic dipole moments and
    flavor-violating processes, receive contributions proportional to the fermion masses, so that
    their running cannot be captured without mass insertions.

    \item We computed the complete one-loop anomalous dimensions of ALP EFTs with both CP-odd and
    CP-even interactions, above and below the electroweak scale, reproducing and extending previous
    results in the literature. These include the dimension-six operators generated once the ALP is
    integrated out, such as the Weinberg operator and the (chromo-)electric and (chromo-)magnetic
    dipole operators. In particular, the on-shell approach makes manifest the relation between operators with
    opposite CP properties, which is hidden in the diagrammatic computation that we performed in
    parallel for comparison. Finally, we evaluated the phase-space integrals also through the
    double-cut method derived from Stokes' theorem, which proved to be technically simpler and more
    efficient.

    \item We classified the physical operators up to dimension six of the most general effective
    gauge theory, with any number of scalar and fermion fields in arbitrary, possibly reducible,
    representations of an arbitrary gauge group, and computed the complete one-loop
    anomalous dimensions of its bosonic and fermionic sectors at order $1/\Lambda^2$. Therefore, the
    renormalization-group equations of any specific model follow from these results through group-theoretic
    manipulations alone, as we verified for the SMEFT, the Low-Energy Effective Field Theory, and
    the scalar $O(n)$ model. 

    \item We showed that helicity remains an active constraint beyond one loop. Each operator is
    assigned the weights $(w,\overline w)=(\ell-h,\ell+h)$, where the length $\ell$ and $h$ are the
    number of particles and the total helicity of its minimal configuration. At $L$ loops, the
    mixing $\mathcal O_j\to\mathcal O_i$ vanishes when $\ell_j-\ell_i=L-1$ and either
    $w_i<w_j-2(L-1)$ or $\overline w_i<\overline w_j-2(L-1)$, in the absence of non-holomorphic
    Yukawa couplings. These zeros follow entirely from tree-level data and are scheme independent
    under finite renormalizations that preserve the operator-length selection rule. Instead, for
    $\ell_i\ge\ell_j$ at two loops, we constructed a scheme in which the mixing vanishes if
    $w_i<w_j-4$ or $\overline w_i<\overline w_j-4$, while, in any other scheme, the same entries
    carry no independent two-loop information, being fixed by one-loop data. Moreover, for operators
    of dimension five through eight, these theorems uncover broad classes of previously unknown
    zeros, from two up to five loops.
\end{itemize}

\subsubsection*{Theoretical Constraints}

A complementary direction concerns the region of parameter space allowed by first principles.
Indeed, properties of the $S$-matrix, such as unitarity, causality, and locality, constrain
the Wilson coefficients irrespective of any experimental information. Among these, partial-wave
unitarity bounds have historically played a central role, since they led to the no-lose Higgs
theorem, which set a target for the LHC. However, their standard formulation applies only to $2\to2$
processes, whereas $2\to N$ amplitudes can grow faster with the energy and are expected to dominate at the high energies of future facilities, and it becomes impractical
for spin-two and higher-spin particles, relevant for gravitational EFTs. Moreover, positivity
bounds restrict the signs of combinations of Wilson
coefficients and, therefore, complement unitarity bounds, which restrict their magnitudes. Along these
lines, this thesis provides the following contributions.
\begin{itemize}
    \item We generalized partial-wave unitarity bounds to processes with arbitrary multiplicity and
    external spins, constructing the angular-momentum basis through the Pauli--Lubanski operator. In
    particular, we provided the basis for $2\to3$ amplitudes at generic $J$, extending previous
    results valid only for $J=0$, and showed that the resulting bounds on dimension-six and
    dimension-eight SMEFT operators supersede the $2\to2$ ones at high energies. Moreover, we
    derived the bounds on the anomalous quartic interactions of photons and gravitons, and analyzed their interplay with positivity bounds.

    \item We derived the partial-wave unitarity bounds on ALP EFTs up to dimension eight, considering also
    weak-violating couplings to leptons, which exhibit an energy enhancement in charged-meson and
    $W$-boson decays. In particular, we highlighted the role of coupled-channel analyses in obtaining the
    tightest bounds. Moreover, we derived the positivity bounds on the dimension-eight operators. Finally, we showed that the size of
    physical effects in several observables is substantially reduced once perturbative unitarity is
    assumed to hold at a given energy scale, and that the bounds on weak-violating couplings are
    typically stronger than current limits from rare decays already at $\sqrt s=1$~TeV.

    \item We provided the first complete implementation of perturbative unitarity
    for the baryon-number-conserving dimension-six SMEFT, covering the full operator basis through
    generic $N\to M$ amplitudes. Moreover, we derived bounds on four-fermion operators that follow from spinning sum rules. We found that these
    constraints are already competitive with, or even stronger than, the experimental bounds above a
    few TeV, especially for four-fermion operators under realistic flavor assumptions, where the sum rules strengthen them further.
\end{itemize}

% \bigskip

% To summarize, the two lines of research provide complementary information on the same Wilson
% coefficients. Indeed, anomalous dimensions describe how they evolve with the energy, while unitarity
% and positivity bounds constrain the region they may populate, and both follow from the same on-shell
% data. 

\subsubsection*{Future Directions}

Several directions follow naturally from these results. First, new-physics scenarios defined above
the TeV scale affect observables measured at or below the GeV scale, such as flavor-violating
processes and electric and magnetic dipole moments, and such a large separation of scales requires
the running at two loops to obtain reliable predictions. In this respect, the master formulae
derived in this thesis already include the most general operator mixing at two loops, where the
anomalous dimensions follow from two-particle cuts of one-loop amplitudes and form factors and from
three-particle cuts of tree-level ones, while the one-loop anomalous dimensions and
$\beta$-functions computed here provide part of the required inputs. Therefore, a natural goal is
the complete two-loop renormalization of the most general effective gauge theory. However, the
one-loop amplitudes and form factors entering the cuts contain rational terms that are not captured
by four-dimensional unitarity cuts, and evanescent operators and
$\gamma_5$ require a consistent choice of renormalization scheme. In this context, the selection
rules derived here reduce the number of independent entries to be computed and provide non-trivial
checks on the result. Moreover, they call for further extensions, such as the holomorphic scheme
beyond two loops, non-linear mixing from multiple operator insertions, and additional zeros based on
angular momentum, which may emerge by decomposing the three-particle cuts into the angular-momentum
basis developed for the unitarity bounds.

Second, on-shell methods have recently been extended to the matching of UV models onto their EFTs,
where the low-energy amplitudes are reconstructed directly in the physical basis from residues and
unitarity cuts of the UV amplitudes~\cite{DeAngelis:2023bmd,Li:2023pfw,Chala:2024llp,Dong:2025wvf}.
Combined with the results of this thesis, they provide a fully on-shell path from a UV model to
low-energy observables, and may eventually allow running and matching to be treated at once, in the
spirit of Ref.~\cite{Giudice:1997ni}. Moreover, this viewpoint is well suited to explain magic
zeros, namely Wilson coefficients that vanish at one-loop matching without an apparent symmetry
reason, such as the leading contribution to the muon dipole moments from weak doublet and singlet
vector-like leptons~\cite{Arkani-Hamed:2021xlp,Craig:2021ksw}. Indeed, in the on-shell approach
these finite contributions follow from phase-space integrals of products of tree-level amplitudes,
which vanish when the integrand is odd under a parity transformation of the phase-space
variables~\cite{DelleRose:2022ygn}.

Third, unitarity bounds, positivity bounds, and spinning sum rules are part of a broader web of
positivity bounds that has emerged from recent advances, including the
EFT-hedron~\cite{Arkani-Hamed:2020blm}, moment bounds~\cite{Bellazzini:2020cot}, and null
constraints~\cite{Caron-Huot:2020cmc,Caron-Huot:2021rmr,Tolley:2020gtv}, which together reveal a
rich network of correlated inequalities on the Wilson coefficients. These constraints can sharpen or
replace model-dependent assumptions in several contexts. Moreover, through a convex-geometry
perspective~\cite{Zhang:2020jyn}, they may even identify the UV completion, among infinitely many
possibilities, once a subset of Wilson coefficients is measured.

Finally, the results of this thesis can be made directly usable in phenomenological studies. Indeed,
an automated implementation of the group-theoretic manipulations would provide the
renormalization-group equations of any specific model from the general ones. Moreover, unitarity and
positivity bounds can be incorporated into statistical analyses, for instance as priors that
restrict the parameter space to physically allowed regions, and as theoretical inputs in event
generation, where they can influence expected distributions and the interpretation of experimental
searches.